% Historical paired-audit draft. Current algorithm-led narrative: draft_v7.tex.
% Standalone: no external figures or bibliography files.
% Replaces v5's unsupported cancellation-exclusive and joint-tail-probability claims.
\documentclass[10pt,a4paper]{article}
\usepackage{kotex}
\usepackage{fontspec}
\IfFontExistsTF{Noto Sans CJK KR}{\setmainhangulfont{Noto Sans CJK KR}}{\setmainhangulfont{Malgun Gothic}}
\usepackage[margin=22mm]{geometry}
\usepackage{amsmath,amssymb,amsthm,booktabs,tabularx}
\usepackage[colorlinks=true,linkcolor=blue,citecolor=blue,urlcolor=blue]{hyperref}
\newtheorem{proposition}{명제}
\newcommand{\norm}[1]{\left\lVert#1\right\rVert}
\title{연합 언러닝 잠복 백도어의 요청 조건부 감사:\\공동 삭제 검증과 식별 한계}
\author{연구 초안}
\date{v6 / 2026년 10월 13일}
\begin{document}
\maketitle

\begin{abstract}
연합 언러닝에서 학습 중 억제된 백도어가 삭제 요청 후 드러날 수 있다.
이때 삭제 후 출력의 큰 변화나 클라이언트 간 반대 방향의 기여만으로 악의를 판정하면,
정상 데이터 이질성 및 비선형 삭제 오차를 공격으로 오인할 수 있다.
본 연구는 기존 학습 중 감시(CCD), 요청 이상 검사(DOSE), 상대 기여 추정(TRACE)을 유지하면서,
요청자만 삭제한 모델과 추정 상대까지 함께 삭제한 모델을 별도 입력에서 비교하는
요청 조건부 공동삭제 검증을 제안한다. 네 개의 실제 개입 모델로 상쇄와 복원을 측정하고,
선형 설명의 비선형 오차 및 악의의 비식별성을 명시한다.
추가 검증은 기존 위험 판정을 취소하지 않는 방식으로 통합한다.
저장된 FUBA, BadFU 변형, 정상 학습 궤적 각 한 개에서 42개 요청--삭제비율 조합을 평가했다.
FUBA에서는 5\%와 10\% 부분 삭제의 변화를 확인했으나 상대 후보에 정상 클라이언트가 포함되었다.
BadFU 변형에서는 공동삭제 구조를 확인하지 못했고 기존 DOSE만 요청 뷰에서 경보를 냈다.
따라서 본 결과는 감사 프로토콜과 제한된 추가 증거를 뒷받침하며,
보편적 공격 탐지, 정확한 공격자 식별, 또는 망각 인증을 입증하지 않는다.
\end{abstract}

\section{문제와 핵심 주장}
FUBA와 BadFU는 학습된 모델의 잠복 성분이 삭제 이후 활성화되는 위협을 다룬다
\cite{fuba,badfu}. 그러나 두 공격의 삭제 단위는 같지 않다.
클라이언트 간 역할 분리가 있는 FUBA와 달리 원래 BadFU는 같은 악성 클라이언트의
backdoor 표본과 camouflage 표본을 사용한다. 공통 방어를 위해서는 공격 이름보다
\emph{어떤 삭제 단위를 어떤 연산자로 제거하고 어디에서 변화를 관측하는가}를 명시해야 한다.

본 연구의 핵심 주장은 다음과 같다.
\begin{quote}
잠복 백도어에 대한 삭제 감사를 요청 조건부 모델 개입으로 정식화하고,
기존 검사의 위험 신호를 유지하면서 추정 상쇄 관계를 별도 입력과 실제 공동 삭제로 확인한다.
이 검증은 출력 변화의 설명을 강화하지만 악의나 삭제 완료를 단독으로 증명하지 않는다.
\end{quote}

추가한 기여 후보는 (i) 삭제 단위를 추상화한 누적 요청 감사,
(ii) TRACE의 후보 선택과 실제 공동 삭제 확인을 분리한 절차,
(iii) 비선형 설명 오차와 정상 반대 기여의 식별 한계 분석이다.
순위 보정, 최소제곱, 반대 정렬 자체는 새로운 이론으로 주장하지 않는다.
새 절차가 기존 방법보다 일반적으로 우수하다는 주장은 다중 시드와 동일 오탐률 비교 이후에 판단한다.

\section{관련 연구와 재사용 범위}
\paragraph{개별 언러닝 기반 탐지.}
MASA는 개별 로컬 모델의 언러닝 행동을 이용해 백도어 모델을 탐지하며,
비동일 분포의 영향을 완화하는 모델 융합을 포함한다\cite{masa}.
이를 단순한 악성 삭제 요청 탐지기로 기술하거나 원리적으로 무력하다고 단정할 수 없다.
UnlearnGuard는 언러닝 과정에서 제출되는 악성 업데이트를 대상으로,
업데이트 추정과 필터링을 이용한 재구성을 제안한다\cite{unlearnguard}.
본 연구의 저장된 잠복 성분 제거 위협과의 차이를 명시한 뒤 동일 조건에서 비교해야 한다.

\paragraph{반사실적 관측.}
BBCaL은 입력에 대한 반사실적 변형과 예측 일관성으로 추론 단계의 백도어 표본을 탐지한다
\cite{bbcal}. 따라서 반사실적 백도어 방어라는 개념 자체를 최초 기여로 주장하지 않는다.
여기서는 입력 변형이 아니라 삭제 단위에 대한 모델 개입과 선택 후 검증을 다룬다.

\paragraph{저장소의 기존 구성요소.}
CCD는 학습 중 누적 기능 기여를 감시하고, DOSE는 동일 삭제량의 비교군을 이용한 요청 검사를,
TRACE는 제한된 최소제곱으로 상대 기여 후보를 찾는 역할을 유지한다.
새 구현은 기존 TRACE 전진 선택 및 box-NNLS 함수를 재사용한다.
FedEraser\cite{federaser}와 재학습은 복구 또는 개입 연산자의 후보다.
이번 실험의 개입은 근사 감산이며, FedEraser까지 포함한 일반적 성능 비교가 아니다.

\section{모형과 알고리즘}
\subsection{삭제 단위, 불변 기준, 관측}
서로 겹치지 않는 삭제 단위를 $\mathcal U=\{1,\dots,K\}$라 하고
누적 삭제 비율을 $r\in[0,1]^K$라 한다. 단위는 클라이언트 또는 데이터 묶음일 수 있다.
학습 기록 $\mathcal H$와 고정된 난수 규약 $\xi$ 아래
\begin{equation}
 M(r)=\mathcal R(\theta_0,\mathcal H,r;\xi),\qquad M(0)=\theta_0
\end{equation}
를 정의한다. 현재 구현은 $M(r)=\theta_0-\sum_jr_jU_j$이며,
누적 기여 $U_j$의 참여율과 집계 가중치를 반영한다.
이 식은 후속 학습의 의존성을 정확히 제거하지 않으므로 망각 인증이 아니다.
요청을 여러 번 나누어 제출해도 $r_t=r_{t-1}+\delta_t$로 누적하며,
부분 제거된 모델로 기준을 바꾸지 않는다.

공개 뷰 $X_p$와 정당하게 접근 가능한 삭제 대상 뷰 $X_f$를 별도로 평가한다.
공격 표적, 역할, 트리거 정답은 탐지에 사용하지 않는다.
삭제 대상 입력이 실제 학습 기록과 일치하는지 확인하는 약정 검증은 별도 전제다.
클라이언트 기록만으로 같은 클라이언트 내부의 임의 데이터 묶음을 제거할 수 있다고 가정하지 않는다.

\subsection{네 모델과 기능 효과}
$S=\operatorname{supp}(r)$, 상대 후보 비율 $a$에 대해
$\operatorname{supp}(a)\cap S=\varnothing$, $a\in[0,1]^K$라 하자.
\begin{equation}
 M_0=M(0),\quad M_S=M(r),\quad M_A=M(a),\quad M_{SA}=M(r+a).
\end{equation}
모든 모델은 같은 기준과 같은 삭제 연산자로 만든다. 예측 확률 $p$의 표현은
\begin{equation}
 \psi_c(p)=\log\max(p_c,\epsilon)-\frac1C\sum_d\log\max(p_d,\epsilon),
 \quad\epsilon=10^{-7}
\end{equation}
이다. 클리핑이 작동하지 않으면 중심화한 로짓과 같다.
뷰 $v$의 입력별 기능 효과는
\begin{equation}
 D_B^v=[\psi(p_{M_B}(x))-\psi(p_{M_0}(x))]_{x\in X_v},\quad B\in\{S,A,SA\}
\end{equation}
이다. 입력 및 클래스 차원을 유지하여 평균에 의해 지워지는 국소 변화를 보존한다.

\subsection{선택과 실제 공동 삭제 확인}
각 뷰를 선택용 $X_v^d$와 확인용 $X_v^h$로 나눈다.
$\rho=\max_jr_j$이고 $E_j=D^d(\rho e_j)$일 때 기존 TRACE로
\begin{equation}
 (\hat A,\hat w)\approx\arg\min_{|A|\le m,\,0\le w_j\le1/\rho}
 \norm{D_S^d+\sum_{j\in A}w_jE_j}^2,\qquad \hat a_j=\rho\hat w_j
\end{equation}
를 계산한다. 뷰별 원소 수의 제곱근으로 정규화한 후 연결한다.
기본 후보 예산 $m=5$는 공격자 수 가정이 아니며 전진 선택의 전역 최적성은 보장하지 않는다.
선택한 집합과 가중치를 고정한 후 실제 $M_A,M_{SA}$를 계산한다.

확인 입력에서 $E_B=\norm{D_B^h}^2$라 할 때, 노름 제곱은 원소별 제곱의 평균으로 정의한다.
\begin{equation}\label{eq:pair}
 C_{S,A}=\frac{E_S+E_A-E_{SA}}{E_S+E_A},\qquad
 R_{S,A}=1-\frac{E_{SA}}{E_S}.
\end{equation}
분모는 구현에서 $10^{-12}$ 이상으로 두고 $E_S>10^{-6}$을 요구한다.
$C$는 요청과 상대의 에너지 대비 공동 잔차를, $R$은 요청 단독 변화 대비 복원을 측정한다.
큰 값은 출력 복원이며 백도어 소멸과 동일하지 않다.

\subsection{위험 판정과 전체 과정}
동일한 전체 삭제 비율과 단위 수의 비교 요청들에서 에너지 중앙값 $\mu_v$와 MAD 척도 $s_v$를 구한다.
이는 데이터 크기와 참여량의 자동 일치를 보장하지 않는다.
\begin{equation}
 z_v=\frac{E_S^v-\mu_v}{\max(s_v,10^{-6})},\qquad
 T=\max_v [z_v]_+[C_v]_0^1[R_v]_0^1.
\end{equation}
기본 추가 경보는 어느 뷰에서 $z_v\ge3$, $C_v,R_v\ge0.5$이고
에너지 및 비어 있지 않은 상대 조건을 만족할 때다. 공개 뷰 TV가 0.3을 초과하면 판정 불가로 둔다.
이 문턱은 휴리스틱이며 $z_v$를 정규분포 유의확률로 해석하지 않는다.

전체 과정은 다음과 같다.
\begin{enumerate}
\item CCD로 학습 중 감사 우선순위를 기록한다. 무경보가 요청 검사 면제를 뜻하지 않는다.
\item 누적 요청에 대해 작은 삭제 단계부터 기존 DOSE와 선택적 공동삭제 검증을 수행한다.
\item 기존 DOSE 위험, 새 검증 위험, 보존 예산 위반 중 하나라도 있으면 보류한다.
      공동삭제 미확인은 기존 경보를 취소할 수 없다.
\item 후보 관계를 증거로 남기고, 요청 전체를 반영한 복구 모델을 별도로 구성한다.
      후보 집합을 정답 공격자로 간주하여 자동 배제하지 않는다.
\item 별도 보존·잔여 공격·망각 검증 후에만 모델을 교체한다.
      재학습 이후에는 새 기준과 기록으로 후속 요청과 재활성화를 감시한다.
\end{enumerate}
현 체크포인트 어댑터는 실제 복구와 독립 망각 검증을 제공하지 않으므로
삭제 완료나 복구 완료 상태를 출력하지 않는다.
후속 학습의 지연 재활성화는 별도 실험이 필요하다.

\section{분석}
\begin{proposition}[가법 효과에서의 반대 정렬]
$D_{SA}=D_S+D_A$이고 $E_S+E_A>0$이면
\begin{equation}
 C_{S,A}=\frac{-2\langle D_S,D_A\rangle}{E_S+E_A}\in[-1,1].
\end{equation}
같은 크기의 완전 반대 효과에서는 $C=R=1$이다.
\end{proposition}
\begin{proof}
$\norm{D_S+D_A}^2=E_S+E_A+2\langle D_S,D_A\rangle$를 식~\eqref{eq:pair}에 대입한다.
Cauchy--Schwarz와 $2uv\le u^2+v^2$로 범위를 얻는다.
\end{proof}

\begin{proposition}[비선형 공동 삭제 오차]
$q=D_{SA}-(D_S+D_A)$이면
\begin{equation}
 \left|C_{S,A}+\frac{2\langle D_S,D_A\rangle}{E_S+E_A}\right|
 \le\frac{2\norm{D_S+D_A}\norm q+\norm q^2}{E_S+E_A}.
\end{equation}
\end{proposition}
\begin{proof}
$\norm{D_S+D_A+q}^2$를 전개하여 남는
$2\langle D_S+D_A,q\rangle+\norm q^2$에 삼각부등식을 적용한다.
\end{proof}
따라서 선택용 선형 잔차가 작다는 것만으로 실제 공동 삭제의 복원을 추론할 수 없다.
일반 비선형 효과에서는 $C$의 하한 $-1$도 성립하지 않는다.
한편 혼합 차분 $q$만으로 탐지하면 가법적인 완전 상쇄($q=0$)를 놓칠 수 있다.

\begin{proposition}[상쇄 관측만으로 악의를 식별할 수 없음]
관측 가능한 모든 입력과 네 개입 모델의 예측 분포가 같은 정상 세계와 공격 세계가 존재하면,
그 관측만을 사용하는 판정기는 두 세계를 항상 구별할 수 없다.
\end{proposition}
\begin{proof}
정상 세계에는 충돌하는 두 학습 목적을, 공격 세계에는 주입과 억제 목적을 배정하되
기능 효과를 같게 구성한다. 관측 분포가 같으므로 어떤 관측 함수의 출력 분포도 같다.
\end{proof}
이 반례는 ``정상 학습에서는 상쇄가 없다''는 전제를 배제한다.
공개 입력이 트리거 영역의 변화를 관측하지 못하는 경우도 탐지 보장의 범위 밖이다.

\paragraph{정상 에피소드 보정의 조건.}
전체 선택·확인·뷰/삭제단계 집계를 고정한 점수 $T$와 정상 에피소드 점수 $T_1,\dots,T_n$에 대해
\begin{equation}
 p=\frac{1+\sum_i\mathbf1\{T_i\ge T\}}{n+1}
\end{equation}
를 사용할 수 있다. 정상 시험/보정 점수가 교환 가능하면 순위 대칭성으로
$\Pr(p\le\alpha)\le\alpha$이다\cite{exchangeability}.
이는 알려진 순위 검정의 적용이며 최소 p값은 $1/(n+1)$이다.
같은 모델의 여러 입력, 동료 클라이언트, 삭제 비율을 독립 학습 에피소드로 세면 안 된다.
현재 실제 실험에는 정상 에피소드 보정 집합이 없으므로 유의수준 보장은 주장하지 않는다.
적응적 요청의 반복 검정에는 추가 조건부 유효성이 필요하다.

\paragraph{비용.}
한 단위 요청에서 예측을 재사용하면 최대 $K+3$개 개입 모델의 순전파가 필요하다.
다중 요청은 비교 요청 조합만큼 추가된다. 모델 상태 생성과 로딩 비용이 존재한다.
역전파가 없다는 설명은 근사 감산 연산자에만 해당하며 재학습 연산자에는 적용되지 않는다.

\section{새로 수행한 실험}
\subsection{프로토콜과 근거 수준}
결과 장부 R-020--R-027에 기록한 로컬 저장 궤적 세 개 및 합성 에피소드를 사용했다. FUBA는 MNIST/LeNet, $K=20$, 100라운드,
참여율 0.4의 IID 한 시드이고 요청자는 억제자 4 및 정상 5,19이다.
BadFU는 CIFAR-10/ResNet-18, 15라운드, 6클라이언트 한 시드의 저장소 변형이며,
요청자는 camouflage 클라이언트 5 및 정상 1,2이다.
이 변형은 원래 논문의 같은 클라이언트 내 표본 삭제 재현이 아니다.
정상 대조군은 MNIST/LeNet, $K=8$, 8라운드, 공격 없는 한 궤적의 전체 8명이다.
총 $(3+3+8)\times3=42$개 요청--삭제비율 조합을 측정했다.

비율은 $0.05,0.1,0.2$, 각 뷰 입력은 256개로 고정했고 128개씩 선택/확인에 사용했다.
공개 입력은 시험 집합 앞 256개, 요청 입력은 해당 클라이언트의 훈련 분할에서
고정 난수로 균등 표집했다. 정확도는 그 다음 1,000개 입력에서 평가했다.
FUBA ASR은 같은 평가 구간의 비표적 899개에 저장 공격 생성기를 적용하여 계산했다.
BadFU ASR은 공개 입력 구간을 제외한 저장 백도어 시험 표본 1,000개에서 계산했다.
역할과 트리거는 개입 선택을 저장한 뒤 채점에만 사용했다.
평가용 GPU 비결정성으로 BadFU 정확도의 일부 재측정에 0.1\%p 차이가 있었다.

표의 반복 비율은 독립 시드가 아니다. 비교군의 규모와 학습 길이도 서로 다르므로
정상 대조군 결과를 FUBA 논문 설정 전체의 오탐률로 일반화하지 않는다.
새로 모델을 학습한 실험이 아니라 저장 궤적에서 실행한 감사 및 평가 실험이다.

\subsection{추가 신호와 한계}
\begin{table}[ht]\centering\small
\caption{새 실제 감사의 경보 수. 분모는 요청--비율 조합으로 서로 상관되어 있다.}
\begin{tabular}{lrrrr}\toprule
대상 & 에너지 단독 & 기존 DOSE & 공동 검증 & DOSE 또는 공동 검증\\\midrule
FUBA 억제자 (3) & 3 & 1 & 2 & 3\\
FUBA 정상 (6) & 0 & 0 & 0 & 0\\
BadFU 변형 요청자 (3) & 0 & 1 & 0 & 1\\
BadFU 변형 정상 (6) & 6 & 0 & 0 & 0\\
공격 없는 정상 (24) & 3 & 0 & 0 & 0\\\bottomrule
\end{tabular}
\end{table}
에너지 단독은 같은 $z\ge3$ 및 최소 에너지 조건을 사용한 소거 실험이다.
공동 확인은 이 실험에서 에너지 단독의 정상 경보 9개를 남기지 않았다.
그러나 기존 DOSE에 대한 정상 오탐 개선은 관측되지 않았다(둘 다 정상 경보 0).
FUBA에서는 기존 flow 기반 검사가 경보를 내지 않은 5\%, 10\%에서 추가 신호가 있었다.
20\%에서는 공개 출력 붕괴 때문에 공동 검증은 판정 불가였고 DOSE가 경보를 냈다.
BadFU에서는 상대 후보를 찾지 못했으며, 20\%에서 요청 뷰의 DOSE만 경보를 냈다.
따라서 두 검사의 AND 결합은 부적절하고 OR 결합의 오탐률은 별도 검증해야 한다.

\begin{table}[ht]\centering\small
\caption{FUBA 억제자 요청의 실제 개입 모델. ACC/ASR(\%). 배포되거나 삭제 완료로 승인된 모델이 아니다.}
\begin{tabular}{lrrr}\toprule
삭제 비율 & 요청만 삭제 ACC/ASR & 공동 삭제 ACC/ASR & 확인 뷰 복원 $R$\\\midrule
0 (기준) & 98.0 / 9.8 & --- & ---\\
0.05 & 98.0 / 27.0 & 98.1 / 8.7 & 0.750\\
0.10 & 96.9 / 49.8 & 98.1 / 9.8 & 0.768\\
0.20 & 44.6 / 98.8 & 96.7 / 11.9 & 0.705\\\bottomrule
\end{tabular}
\end{table}
공동 삭제는 ASR을 기준 모델 근처로 되돌렸으나 백도어 소멸을 입증하지 못했다.
10\%에서 선택된 상대는 $\{0,3,17,13,15\}$로 실제 주입자 $\{0,1,2,3\}$와 일치하지 않는다.
정상 후보가 섞였으므로 자동 배제 근거로 사용할 수 없다.
별도로 저장되어 있던 정답 집합 제외 재학습 모델은 같은 평가에서 ACC 98.7\%, ASR 0.44\%였다.
이는 이미 알려진 제외 집합을 쓴 참조이며 새 후보 집합으로 재학습한 결과가 아니다.

BadFU 변형은 삭제 전 ACC 83.8\%, ASR 26.2\%로 충분히 강한 잠복 재현이라고 보기 어렵다.
요청자의 20\% 삭제 후 ACC 82.5\%, ASR 38.0\%였고 공동 검증은 개선하지 못했다.
이 음성 결과를 원래 BadFU 방어 성공으로 집계하지 않는다.

\subsection{메커니즘 소거와 반례}
선형 상쇄, 공개 뷰에 보이지 않는 상쇄, 비선형 공동 개입 실패,
선택 입력에만 맞는 상대, 정상 반대 기여, 작은 정상 효과를 각 50회 생성했다.
실제 공동 삭제 확인은 비선형 실패 50/50을, 별도 입력 확인은 선택 과적합 50/50을 거부했다.
선택용 선형 복원만 사용하면 두 경우 모두 50/50에서 복원을 잘못 예상했다.
반면 정상 반대 기여는 공격 상쇄와 동일한 관측을 만들었으므로 50/50에서 추가 경보가 났다.
이는 상쇄의 비식별성에 대한 구성적 반례이며 현실 오탐률 추정이 아니다.

별도 합성 정상 에피소드 199개로 고정 파이프라인 점수를 보정하고
새 정상 에피소드 200개를 평가했을 때 $\alpha=0.05$에서 4개가 기각되었다.
이 결과는 해당 합성 생성기에만 해당하며 실제 FL 정상 보정 집합을 대체하지 않는다.
회복/분할 요청/파일 가중치/새 판정 통합 등을 포함한 동작 검사 37개가 통과했다.

\subsection{기존 결과의 위치}
기존 CCD의 FUBA 논문 설정 6개 캐시 재생은 저장된 첫 경보 라운드
IID 13/15/11, non-IID 16/15/11 및 후보 집합과 일치했다.
이는 기존 구성요소의 재사용 회귀 확인이며 새 공동삭제 검증의 6시드 성능이 아니다.
기존 복구 결과와 FedEraser 결과도 별도 프로토콜의 기록으로 남기고,
새 통합 방법의 종단간 결과로 합산하지 않는다.

\section{후속 검증과 결론}
출판 수준의 성능 주장을 위해서는 다음 검증이 필요하다.
첫째, 같은 학습 설정의 공격 없는 IID/Dirichlet/클래스 소유자 궤적으로
문턱 보정과 시험을 분리하고, 기존 DOSE 및 에너지 단독과 동일 오탐률에서 비교해야 한다.
둘째, 원래 BadFU의 같은 클라이언트 내 묶음 삭제와 충분히 잠복한 모델을 재현해야 한다.
셋째, 동일 정보·연산자·비용 조건에서 MASA, UnlearnGuard, FedEraser, 전체 재학습을 비교해야 한다.
넷째, 추정 후보로 실제 복구한 뒤 보존·망각·잔여 ASR·후속 학습 재활성화를 측정해야 한다.
다섯째, 문턱·뷰·후보 예산을 고정한 뒤 학습 중 방어를 고려하는 적응형 공격을 평가해야 한다.
저장 업데이트의 단순 배율 조정은 마지막 실험을 대체하지 않는다.

본 연구는 기존 방법을 유지하며 상쇄 설명에 실제 개입과 별도 입력 검증을 추가했다.
제한된 실험에서 새 검증은 일부 FUBA 요청에 추가 증거를 제공하고 단순 에너지 오경보를 걸렀지만,
정상 반대 기여와 BadFU 변형의 실패는 보편적 탐지 주장을 배제한다.
따라서 현 단계의 기여는 검증 가능한 감사 절차와 그 한계이며,
정확한 공격자 식별 및 안전한 언러닝 완료는 독립된 실증 과제로 남는다.

\begin{thebibliography}{9}
\bibitem{fuba} FUBA: Backdoor Federated Learning via Federated Unlearning.
\href{https://ieeexplore.ieee.org/document/11231135/}{IEEE Transactions on Artificial Intelligence, DOI:10.1109/TAI.2025.3630110}.
\bibitem{badfu} B. Lu et al. BadFU: Backdoor Federated Learning through Adversarial Machine Unlearning.
\href{https://arxiv.org/abs/2508.15541}{arXiv:2508.15541}, 2025.
\bibitem{masa} J. Xu, Z. Zhang, R. Hu. Identify Backdoored Model in Federated Learning via Individual Unlearning.
\href{https://openaccess.thecvf.com/content/WACV2025/html/Xu_Identify_Backdoored_Model_in_Federated_Learning_via_Individual_Unlearning_WACV_2025_paper.html}{WACV}, 7949--7958, 2025.
\bibitem{unlearnguard} W. Wang et al. Poisoning Attacks and Defenses to Federated Unlearning.
\href{https://arxiv.org/abs/2501.17396}{WWW / arXiv:2501.17396}, 2025.
\bibitem{bbcal} BBCaL: Black-box Backdoor Detection under the Causality Lens.
\href{https://openreview.net/notes/edits/attachment?id=cJ3e8RwIZu\&name=pdf}{공개 심사 원고}; TMLR, 2024.
\bibitem{federaser} G. Liu et al. FedEraser: Enabling Efficient Client-Level Data Removal from Federated Learning Models. IWQoS, 2021.
\bibitem{exchangeability} A. K. Kuchibhotla. Exchangeability, Conformal Prediction, and Rank Tests.
\href{https://arxiv.org/abs/2005.06095}{arXiv:2005.06095}, 2020/2021.
\end{thebibliography}
\end{document}
